% Question 3 -- draft so far. 3.(a): definitions and daily problem are written (they only restate
% the brief); the properties (problem class, convexity, feasibility) are left as TODO with hints.
%
% NOTE ON 2.(c): Question 3 starts from the quadratic disutility of 2.(c), which is not written
% yet. Below the disutility is kept generic, D_t(l_t), as in 2.(a) (convex, differentiable). This
% is enough to write 3.(b): the KKT conditions only need D_t'(l_t). Once 2.(c) is written, D_t is
% replaced by the quadratic disutility and the properties below are specialised.
%
% 3.(c) (qualitative hypotheses), Ignacio's note: the reasoning -- does the consumer ever
% consume above its hourly reference, and above E^min -- follows from a cost-benefit argument
% (disutility of deviating vs. cost/benefit of importing/exporting), not from duality, so it
% does not need 2.(c) to be implemented. Not written yet.

\subsection{(a) Formulation with a minimum daily energy}

\subsubsection{New input data and constraint}
Starting from the model of Question~2.(c), the consumer must now consume at least a minimum
total energy $\Emin$ (kWh) over the whole day. This is the only new input datum, and it adds a
single new constraint:
\begin{equation}
    \sum_{t \in \T} \load \geq \Emin.
    \label{eq:c_min_energy}
\end{equation}

\subsubsection{Daily problem}
The daily problem keeps the objective of Question~2 and all the constraints of Question~1, and
adds \eqref{eq:c_min_energy}. With $D_t$ the disutility of the deviation of the load from the
reference $\lref$ (in Question~3, the quadratic disutility of 2.(c)):
\begin{align}
    \max_{\{\load, \qimp, \qexp, \pvprod\}_{t \in \T}} \quad
    & \sum_{t \in \T} \Big( - D_t(\load) - \pimp \, \qimp + \pexp \, \qexp - \cpv \, \pvprod \Big)
    \label{eq:q3_objective} \\
    \text{s.t.} \quad
    & \load = \pvprod + \qimp - \qexp, & \forall t \in \T, \nonumber \\
    & \Lmin \leq \load \leq \Lmax, & \forall t \in \T, \nonumber \\
    & 0 \leq \pvprod \leq \PVmax, & \forall t \in \T, \nonumber \\
    & \qimp \geq 0, \quad \qexp \geq 0, & \forall t \in \T, \nonumber \\
    & \sum_{t \in \T} \load \geq \Emin. & \nonumber
\end{align}

\subsubsection{Change in the properties}
\begin{itemize}
    \item \textbf{Decomposability:} constraint \eqref{eq:c_min_energy} is the first constraint
    of the assignment that links different hours, since it involves the loads of all 24
    hours. The daily problem can therefore no longer be split into 24 independent hourly
    problems: it has to be solved as a single 24-hour problem.
    \item \todo{Problem class and convexity compared with Question 2. Hint: look at what the
    new constraint is (linear or not) and at the class of the objective of 2.(c); does adding
    it change the class or the convexity of the problem?}
    \item \todo{Feasibility: give a condition on $\Emin$, the reference profile and the load
    bounds under which the problem is feasible. Hint: think about the largest total energy the
    consumer can consume in one day given the load bounds, and whether the reference profile
    actually restricts what is feasible. This condition is needed again in 3.(b), to justify
    strong duality.}
\end{itemize}

\subsection{(b) Lagrangian, dual problem and KKT conditions of the daily problem}
% HINTS (not final text) for whoever writes 3.(b):
% - Set up as in 1.(b): minimisation, h(x) = 0 and g(x) <= 0, one multiplier per constraint. The
%   hourly multipliers of 1.(b) are reused (macros \lbal, \muLlo, ..., all with the hour index t);
%   the new one is \muE (already in main.tex), a single scalar for the whole day.
% - The Lagrangian is the sum over t of the hourly terms of 1.(b), with D_t in place of the
%   linear utility, plus the term of the minimum-energy constraint.
% - In the stationarity condition of the load, D_t'(l_t) appears and \muE appears in every
%   hour: this is how the multiplier couples the hours in the KKT conditions.
% - Strong duality and necessity/sufficiency of KKT: same route as in 1.(b) (problem class from
%   (a), Slater's condition reduces to feasibility when the constraints are affine), and the
%   feasibility condition of 3.(a) is what is needed here.
Following the brief, the multiplier of the minimum-energy constraint \eqref{eq:c_min_energy}
is denoted by $\muE \geq 0$. It is a multiplier of the same kind as those of Question~1.(b):
it belongs to an inequality, which in the form $g(x) \leq 0$ reads
$\Emin - \sum_{t \in \T} \load \leq 0$, and it is in DKK/kWh. Unlike the hourly multipliers, it
has no hour index, since the constraint applies to the whole day.

\todo{Lagrangian, dual problem and KKT conditions of the daily problem, with the multiplier
of \eqref{eq:c_min_energy}. Does strong duality hold?}
